\documentclass[12pt]{article}

% Layout.
\usepackage[top=1in, bottom=0.75in, left=1in, right=1in, headheight=1in, headsep=6pt]{geometry}

% Fonts.
\usepackage{mathptmx}
\usepackage[scaled=0.86]{helvet}
\renewcommand{\emph}[1]{\textsf{\textbf{#1}}}

% TiKZ.
\usepackage{tikz, pgfplots}
\usetikzlibrary{calc}

% Misc packages.
\usepackage{amsmath,amssymb,latexsym}
\usepackage{graphicx}
\usepackage{array}
\usepackage{xcolor}
\usepackage{multicol}
\usepackage{fancyhdr}
\pagestyle{fancy}

\usepackage[parfill]{parskip}


% Misc.
\renewcommand{\d}{\displaystyle}
\newcommand{\ds}{\displaystyle}
\def\bc{\begin{center}}
\def\ec{\end{center}}
\def\be{\begin{enumerate}}
\def\ee{\end{enumerate}}

\newcommand{\ans}[1][1in]{\rule{#1}{.5pt}}


\newcommand{\corrections}[3][l]{%
  % optional #1 = alignment: l (default), r, c
  % #2 = height (with units, e.g. 1in)
  % #3 = proportion of \linewidth (e.g. .6)
  \par\noindent
  \makebox[\linewidth][#1]{%
    \begin{tikzpicture}[baseline]
      \draw (0,0) rectangle (#3\linewidth,-#2);
      \node[anchor=south west, font=\it\small, inner sep = 0] at (0,1pt) {Write corrections here};
    \end{tikzpicture}%
  }%
  \par\vspace{2pt}
}


\lhead{Math F113X: Quiz 2}

\begin{document}

\vspace{1cm}

\noindent \textbf{Name:} \hrulefill \quad  score:\ans[1cm] \ / 10 \\


There are 10 points possible on this quiz. No aids (book, notes, etc.)
are permitted. You may use a non-programmable calculator.  \emph{Show
all work and supporting calculations for full credit. Explain how you
get your answers.}

%Do not write in the ``write corrections here'' boxes while you're taking the main quiz.


\be
\item (3 points) Consider the weighted voting system $[31: 16, 12, 10, 3]$

\begin{enumerate}
    \item What is the meaning of the number \emph{31} in the weighted voting system?

\vfill

    \item Why is $\{P_1, P_2, P_4\}$ a \emph{winning coalition}? Explain your answer with a calculation and a few words.

\vfill

    \item Which players, if any, have \emph{veto power}? Explain your reasoning.
    
    \vfill

\end{enumerate}

\item (2 points) Consider the weighted voting system $[q: 12, 8, 2, 1]$.

\begin{enumerate} 
    \item What is the \emph{smallest possible quota value}? Explain how you know.
   \vfill

    \item Determine a value of $q$ that results in a \emph{dictator}. Explain which player is the dictator and how you know. 
    
    \vfill
\end{enumerate}

\newpage

\item Consider the weighted voting system $[17: 13, 9, 6, 2]$ and the list of all possible coalitions.

{\renewcommand{\arraystretch}{2.5}
\begin{tabular}{c | c | c || c | c | c}
coalition & weight & winning? & coalition & weight & winning?\\ \hline \hline
$P_1 \ P_2$ & & & 	$P_1\ P_2\ P_3$ 		&  & \\ \hline 
$P_1\ P_3$ & & & 	$P_1\ P_2\ P_4$ 		& 24 & \\ \hline 
$P_1\ P_4$ & & & 	$P_1\ P_3\ P_4$ 		& 21 & \\ \hline
$P_2\ P_3$ & 15 & & 	$P_2\ P_3\ P_4$ 		& 17 & \\ \hline
$P_2\ P_4$ & 11 & & 	$P_1\ P_2\ P_3\ P_4$ 	& 30&  \\ \hline
$P_3\ P_4$ & 8 &  \\
\end{tabular}
}





\begin{enumerate}
    \item (1 point) In the table above, fill in the \emph{missing weights} of each coalition in the appropriate column.
    
    \item (1 point) Put a check $\checkmark$ in the column labeled ``winning?'' to identify all winning coalitions, and mark  $\times$ for non-winning coalitions.

    \item (2 points) In each winning coalition, \emph{underline the critical players}.
    
    \item (1 point) Are there any \emph{dummies}? Explain why or why not, with a sentence or short computation and sentence.

%    \item (1 point) Using this information, compute the Banzhaf Power Index of each player by filling in the following table.

\end{enumerate}

%\bigskip
%
%%{\renewcommand{\arraystretch}{2.5}
%\begin{tabular}{c | p{1in} | c | c }
%Player & &  power (fraction) &  power (\%) \\ \hline
%&&&\\
%&&&\\&&&\\&&&\\&&&\\&&&\\&&&\\&&&\\&&&\\
%\end{tabular}
%%}

\end{enumerate}

\end{document}
%%% Local Variables:
%%% mode: LaTeX
%%% TeX-master: t
%%% End:
